% Pakete laden
\usepackage[utf8]{inputenc} % Zeichenkodierung
\usepackage[ngerman,english]{babel} % Sprache
\usepackage{microtype}      % Schriftbild und den Zeilenumbruch verbessern
\usepackage{graphicx}       % Grafiken
\usepackage{csquotes}       % Zitate
\usepackage[style=ieee]{biblatex} % Literaturverzeichnis
\usepackage{fancyhdr}       % Kopf- und Fußzeilen
\usepackage[T1]{fontenc}    % Schriftart
\usepackage{uarial}         % Arial
\usepackage{parskip}        % Abstand zwischen Absätzen
\usepackage[nopostdot,nonumberlist,acronym,automake,toc,section]{glossaries} % Abkürzungsverzeichnis
\usepackage{titlesec}       % Section- und Subsection-Überschriften
\usepackage[colorlinks=false,pdfborder={0 0 0}]{hyperref} % Hyperlink für Inhaltsverzeichnis und Quellen
\usepackage{nameref}
\usepackage[nottoc,notbib,section]{tocbibind}
\usepackage{url}            % bessere Darstellung von Links
\usepackage{enumitem}       % Aufzählungen
\usepackage{booktabs}       % Tabellen mit schöneren Linien
\usepackage{multirow}       % Tabellen mit mehreren Zeilen
\usepackage{longtable}      % Für mehrseitige Tabellen
\usepackage{makecell}       % Für Zellen in Tabellen
\usepackage{lipsum}         % Dummy Text
\usepackage{xcolor}
\usepackage{xparse}
\usepackage{amssymb}      % Mathematische Symbole wie \mathbb
\usepackage{needspace}
\usepackage{float}


% Ränder einstellen
\usepackage[
    left=2.5cm, 
    right=2.5cm, 
    top=0cm, 
    bottom=2.5cm,
    nomarginpar,
    includehead,
    headheight=2.5cm,
    headsep=1cm
]{geometry}


% Anpassung der Referenzformatierung
\NewDocumentCommand{\fautoref}{m}{\textbf{\autoref{#1}}}
\NewDocumentCommand{\fsecref}{m}{\textbf{\hyperref[#1]{\ref*{#1} \nameref*{#1}}}}
\NewDocumentCommand{\shortsecref}{m}{\textbf{\hyperref[#1]{\nameref*{#1}}}}

% Definiere Bibliografie-Treiber für @video
\DeclareBibliographyAlias{video}{online}
\DeclareBibliographyAlias{software}{online}

% Schriftart ändern
\renewcommand{\familydefault}{\sfdefault}
\usepackage{setspace}

% Setzt den Zeilenabstand
\onehalfspacing

% Setze den Abstand zwischen Absätzen
\setlength{\parskip}{0.75cm}

% Abstand vor und nach Section- und Subsection-Überschriften anpassen
\titlespacing{\section}{0pt}{1.5cm}{0.75cm}
\titlespacing{\subsection}{0pt}{1.25cm}{0.5cm}
\titlespacing{\subsubsection}{0pt}{0.75cm}{0.25cm}


